\subsubsection{Search for Applicant}

In this use case, the Hall Office will be able to conduct searches for applicant data using either the GST ID or the roll number of the applicant. This feature makes it easy for one to be able to get the applicant's data within a short time.Moreover, the system helps display the results in a systematic way, which allows the Hall Office to easily find the right candidate from amongst all those who have applied. The system makes the process of data search much easier and faster for the Hall Office. The search feature also helps minimize any risk of errors as the records can be easily identified using GST ID or roll numbers. Another advantage of the search functionality of the system is that it gives the Hall Office immediate access to all the information about each and every applicant.

\begin{table}[H]
    \centering
    \small
    \renewcommand{\arraystretch}{1.4} 
    \caption{Use Case Table: Search Applicant (UC-3.1)} % ক্যাপশন উপরে আনা হয়েছে
    \label{tab:search-applicant}
    \begin{tabularx}{\textwidth}{|l|X|}
        \hline
        \rowcolor[gray]{0.9} {Description of Characteristics} & {Details} \\ \hline
        Use Case ID & UC-3.1 \\ \hline
        Use Case Name & Search Applicant \\ \hline
        Actors & Hall Office \\ \hline
        Purpose & To locate an applicant’s details in the system. \\ \hline
        Pre-condition & The user must be logged into the system and have the GST Applicant ID or Roll. \\ \hline
        Scenarios & 
        \begin{enumerate}[nosep, leftmargin=*, after=\vspace{-\baselineskip}, before=\vspace{-0.5\baselineskip}]
            \item User navigates to the Hall Office dashboard.
            \item User enters the GST ID or Roll Number.
            \item System fetches and displays the applicant details.
        \end{enumerate} \\ \hline
        Occurrence & Often \\ \hline
    \end{tabularx}
\end{table}
\FloatBarrier